% SYN-ACK WIRE -- user manual for the Coleco Adam, with a FujiNet
% Build: make   (or make in docs/: the rulebook and every manual)
%
% The style, the theme macros and the furniture the shared rules text needs are in ../../style/synack.sty;
% the rules text is ../../rules/*.tex; the Adam's names for things (its keys) are in platform.tex.
% Screens are in screens/: dumps of the Adam client's display, drawn by client/tools/render_adam.py
% (client/tools/capture.sh with PLATFORM=adam makes them).
\documentclass[12pt,letterpaper]{article}

\usepackage{synack}
% The Adam's names for things, for the shared style and text (../../style/synack.sty, ../../rules/).
\renewcommand{\DocKind}{USER MANUAL // ADAM}
\renewcommand{\PlatformShort}{Adam}                        % as in "on the Adam"
\newcommand{\PlatformLong}{Coleco Adam}                    % the title page
\newcommand{\PlatformMachines}{COLECO ADAM \quad // \quad FUJINET ADAM}   % the title page, the screen chapter
\newcommand{\SignalLost}{the title bar turns magenta and reads \code{NO SIGNAL}}  % the link is quiet
\renewcommand{\KeyReturn}{RETURN}                          % the Adam's names for the keys that confirm
\renewcommand{\KeyEsc}{ESCAPE}                             % and back out
\renewcommand{\HostAsk}{HOW MANY SEATS? \Rule{minplayers}-\Rule{maxplayers}}
\renewcommand{\TooLate}{TOO LATE: LAST SEAT TAKEN}

% Where the screen sits in a capture, in image pixels (see \capgeom): client/tools/render_adam.py at --scale 3
% draws all of the TMS9918A's 32 x 24 cells, so a callout's (column, row) is the cell's.
\makeatletter
\newcommand{\rendered}{\capgeom{0}{0}{24}{24}{768}{576}}
\makeatother

\graphicspath{{screens/}}

\makeatletter
% The rules' captures are staged games (server/stage.py, scenarios.toml) on the board of the drawing
% under "1. Spike a node", so the drawing, the text and the screen tell one story.
\expandafter\def\csname shot@spike\endcsname{\screen{spike.png}{CAPTURE \quad \MakeUppercase{\V{spike}} A \MakeUppercase{\V{node}}}%
  {The drawing above, as ZER0COOL sees it. The move hints \code{1}--\code{4} sit on its four targets: they are the
  \Vp{exploit} in your \V{kit} that run there, as the info line lists them. Press the \V{exploit}'s number to run
  it. On the controller a cursor is under the first that can run, a box on its \V{node}: the stick moves it, the
  left fire button runs it.}}
\expandafter\def\csname shot@broker\endcsname{\screen{infiltration.png}{CAPTURE \quad \MakeUppercase{\V{infiltration}}}%
  {Target 2: ZER0COOL ran \node{C6}, next to the lone \V{spike} on \node{B6}. The question line lists, inverse, the
  syndicates with no \V{hive} on the \V{grid}: press one's letter to \V{broker} the new \V{hive} to it (or move the
  cursor under it and press the left fire button).}}
\expandafter\def\csname shot@hivejack\endcsname{\screen{tied-root.png}{CAPTURE \quad \MakeUppercase{\V{hivejack}}: A TIED \MakeUppercase{\V{root}}}%
  {The example below, as ZER0COOL sees it. The spike on \node{G3} joined \Syn{5}{name} (4 \Vp{node}) and \Syn{3}{name}
  (3); the log shows the tied \V{root} paid. The info line says how many of the \V{absorbed} \Vp{coin} you hold:
  type how many to \V{cashout} and press \key{RETURN}, then how many to \V{swap}, an even number, and \key{RETURN}
  again. The stick counts them too.}}
\expandafter\def\csname shot@buy\endcsname{\screen{spread-buy.png}{CAPTURE \quad BUY \MakeUppercase{\V{coin}}}%
  {Target 3: \node{E1} spread \Syn{5}{name} to 5 \Vp{node}, and its price is already up. While you buy, the
  \Vp{hive} panel shows each price and how many the \V{bank} has left. Each syndicate letter you type buys one of
  its \V{coin}, up to \Rule{maxbuy} in any mix; \key{BACKSPACE} takes one back, \key{RETURN} buys. On the
  controller an arrow picks the syndicate and the left fire button adds one.}}
\expandafter\def\csname shot@nullroute\endcsname{\screen{nullroute.png}{CAPTURE \quad \MakeUppercase{\V{lockdown}}: THE QUESTION}%
  {ZER0COOL's \V{spike} on \node{H2} spread \Syn{2}{name} to \Rule{safesize} \Vp{node}: \V{safe}, and with
  \Syn{5}{name} and \Syn{0}{name} already \V{safe}, every \V{hive} is. With too little \Game{currency} for any
  \V{coin}, the prompt is the question alone: \key{Y} calls the \V{nullroute}, \key{N} plays on.}}
\expandafter\def\csname shot@lobby\endcsname{\screen{lobby-join.png}{CAPTURE \quad JOINING A LOBBY}%
  {ZER0COOL hosts and MAMBA is waiting; the Adam has just joined. The host clock started when MAMBA arrived, and
  every newcomer is told how long the host has left.}}
\expandafter\def\csname shot@final\endcsname{\screen{game-over.png}{CAPTURE \quad THE END OF A GAME}%
  {ZER0COOL called it. \VC{root} and \V{sudo} were paid for every \V{hive}, all \V{coin} was sold, and the log ends
  with each \V{player}'s final \Game{currency} and the top \V{player}. The table is a lobby again: the host's prompt
  asks for the seats of a rematch.}}
\makeatother


\begin{document}

% ================================================================== title
\begin{titlepage}
\machine
\vspace*{1.6in}
{\color{amber}\rule{\linewidth}{2pt}}\par\vspace{14pt}
{\fontsize{44}{48}\selectfont \game\par}\vspace{10pt}
{\fontsize{12}{16}\selectfont\color{steel} SEVEN SYNDICATES. ONE REGION. EVERY \MakeUppercase{\V{node}} FOR SALE.\par}
\vspace{14pt}{\color{amber}\rule{\linewidth}{2pt}}\par
\vspace{0.35in}
{\fontsize{24}{30}\selectfont USER MANUAL\par}\vspace{4pt}
{\fontsize{14}{18}\selectfont\color{amber} FOR THE \MakeUppercase{\PlatformLong}\par}
\vspace{0.4in}
{\fontsize{12}{16}\selectfont \PlatformMachines\par
\Rule{minplayers}-\Rule{maxplayers} \MakeUppercase{\Vp{player}} \quad // \quad BOTS FILL EMPTY SEATS\par}
\vfill
{\fontsize{8.5}{11}\selectfont\color{steel}
AN UNOFFICIAL CYBERPUNK RE-SKIN OF ACQUIRE BY SID SACKSON.\par
ACQUIRE IS A TRADEMARK OF HASBRO, INC. NOT AFFILIATED WITH OR ENDORSED BY HASBRO.\par}
\end{titlepage}

% ================================================================== 01
\renewcommand{\secstatus}{THE SETTING \quad // \quad THE OBJECT \quad // \quad THE END}
\chapter{The Network Grid}
\begin{multicols}{2}
\rules{opening}
\end{multicols}

% ================================================================== 02
\renewcommand{\secstatus}{REGION \quad // \quad STATUS \quad // \quad HANDLE}
\chapter{\VC{jackin}}
\game\ runs on a Coleco Adam with a FujiNet. Pick a \game\ server in the FujiNet lobby, which loads the game
itself, or in FujiNet's CONFIG mount \code{syn-ack-wire-adam.ddp} from the host \code{syn-ack-wire.fly.dev} and boot
it. The lobby passes the server you pick through FujiNet's SD card: without one, CHIBA-GRID-7 is the server you join.

The title dials the server as soon as it appears. Type a handle, up to 8 characters, and press \key{RETURN}. From
the FujiNet lobby, the handle starts as your lobby name. The controller dials a handle too: up and down change the
letter the cursor is on, left and right move along the handle, the right fire button deletes a letter, and
\key{\#} on the keypad jacks in.

\begin{capture}{title.png}{CAPTURE 01 \quad TITLE SCREEN \quad LOBBY OPEN}
  \callout{1}{L}{1}{15}
  \callout{2}{R}{22}{16}
  \callout{3}{L}{8}{17}
  \callout{4}{L}{1}{18}
  \callout{5}{R}{17}{19}
  \callout{6}{L}{1}{20}
  \callout{7}{R}{20}{21}
  \callout{8}{L}{1}{22}
  \callout{9}{R}{24}{23}
\end{capture}

\begin{legend}
\li{1}{\textbf{Region.} The server's \V{grid}: one game at a time, up to six seats.}
\li{2}{\textbf{Status.} Whether you can join, and how: \code{LOBBY OPEN}, \code{IN PROGRESS}, \code{TABLE FULL}.
It stays up to date while you type.}
\li{3}{\textbf{Handle.} Your name at the table. The cursor shows inverse: past the end, where a letter goes; or,
with the controller, on the letter it changes.}
\li{4}{\textbf{Keys, and the controller's.} Type the handle, or dial it with the stick.}
\li{5}{\key{RETURN}, or \key{\#}: \VC{jackin}.}
\li{6}{SmartKey \key{I}, or \key{*}, opens the instructions, here or in the middle of a game.}
\li{7}{SmartKey \key{II}, or \key{*} with the left fire button, turns the move hints on or off.}
\li{8}{SmartKey \key{III}, or \key{*} with the right fire button, turns the sound on or off.}
\li{9}{\textbf{The prompt line.} What the screen expects from you next; a refusal, such as a handle already
taken, shows here, on magenta.}
\end{legend}

% ================================================================== 03
\renewcommand{\secstatus}{32 COLUMNS \quad // \quad COLOUR \quad // \quad THE CURSOR}
\chapter{The Screen}
\begin{sloppypar}
Everything you need is on one screen. The Adam has room for 32 columns of letters where the other machines have
40, so \game\ lays its screen out as on the NES: the syndicates by their short names, long log lines over two
rows. Colour does what inverse does elsewhere: the \V{grid} and the seats on green, cash and every other
syndicate's \V{coin} in yellow, and the bottom two lines magenta while the decision is yours. A yellow cursor marks
what the controller's left fire button would answer. This is early in a game, and it's your turn: pick an
\V{exploit} to run.
\end{sloppypar}

\begin{capture}{spike-early.png}{CAPTURE 02 \quad \MakeUppercase{\V{spike}} PROMPT \quad 3 SEATS \quad EARLY GAME}
  \callout{1}{L}{1}{0}
  \callout{2}{L}{2}{6}
  \callout{3}{R}{9}{4}
  \callout{4}{R}{27}{3}
  \callout{5}{R}{21}{10}
  \callout{6}{L}{1}{13}
  \callout{7}{L}{1}{20}
  \callout{8}{L}{1}{22}
  \callout{9}{R}{30}{23}
\end{capture}

\begin{legend}
\li{1}{\textbf{Title bar.} The region and the link. It turns magenta and reads \code{NO SIGNAL} when the server goes quiet.}
\li{2}{\textbf{The \V{grid}.} Inverse letters are \Vp{hive}; \node{\Glyph{spike}} is a lone \V{spike}. A digit marks
a \V{node} an \V{exploit} in your \V{kit} can run on. The last move flashes for a few seconds.}
\li{3}{\textbf{The box.} The \V{node} the \V{exploit} under the cursor would \V{spike}.}
\li{4}{\textbf{\VCp{hive}.} Each syndicate's letter, \V{hive} size and short name. While you buy, its price and how
many the \V{bank} has left.}
\li{5}{\textbf{Last move.} The \V{node} \V{spike}d most recently.}
\li{6}{\textbf{\VCp{player}.} \code{>} marks you, \code{*} a bot, \code{?} someone away. Account, then
\V{coin} held, every other syndicate's in yellow. An inverse name has the turn.}
\li{7}{\textbf{The log.} What happened at the table lately; a long line goes on, indented, on the next row.}
\li{8}{\textbf{Your \V{kit}.} Your \Rule{kitsize} \Vp{exploit} and the \Vp{node} they run on; an inverse one
can't run now. The cursor is under the one the left fire button would run.}
\li{9}{\textbf{The prompt.} What you're asked, with seconds left on the shot clock. The two bottom lines turn
magenta when the decision is yours.}
\end{legend}

% ================================================================== 04
\renewcommand{\secstatus}{THE KEYBOARD \quad // \quad THE CONTROLLER \quad // \quad SMARTKEYS}
\chapter{Keys and Controller}
\game\ is played with the Adam's keyboard, or with the controller in port 1, or both: each question on the prompt
line names the keys it takes, and a yellow cursor sits on the answer the controller's left fire button would give.
The game checks an answer before sending it and says on the prompt line what's wrong, such as
\code{NOT ENOUGH NEW YEN}.

\subsection{Any time}
\begin{keytable}
\kr{\key{I}\ \ \key{*}}{Open the instructions, on the title screen or in the middle of a game: SmartKey \key{I}, or \key{*} on the controller's keypad on its own.}
\kr{\key{II}\ \ \key{*}+\key{L}}{Move hints on or off: the digits on the \V{grid} that show where each \V{exploit} in your \V{kit} can run. On the controller, \key{*} with the left fire button.}
\kr{\key{III}\ \ \key{*}+\key{R}}{Sound on or off. A chord tells you it's on. On the controller, \key{*} with the right fire button.}
\end{keytable}

\subsection{The title screen}
\begin{keytable}
\kr{Letters, digits}{Type your handle. \key{BACKSPACE} deletes the last character.}
\kr{\key{RETURN}\ \ \key{\#}}{\VC{jackin}.}
\kr{The stick}{Up and down change the letter the cursor is on: A--Z, 0--9, then the signs a handle can have. Left and right move along the handle. The right fire button deletes the letter.}
\end{keytable}

\subsection{The instructions}
\begin{keytable}
\kr{\key{RIGHT}}{Next page: the arrow key, or the stick.}
\kr{\key{LEFT}}{Previous page.}
\kr{\key{ESCAPE}\ \ \key{RETURN}}{Back to where you were. SmartKey \key{I}, either fire button and \key{\#} do the same.}
\end{keytable}

\subsection{Hosting}
\begin{keytable}
\kr{\key{\Rule{minplayers}}--\key{\Rule{maxplayers}}}{Start the game with that many seats. Bots take the seats nobody has joined.}
\kr{The stick}{Up and down set the number; the left fire button starts.}
\end{keytable}

\subsection{Your turn}
\begin{keytable}
\kr{\key{1}--\key{6}}{\VC{spike}: run that \V{exploit} in your \V{kit}. An inverse one can't run now; the prompt line says why.}
\kr{Letters}{\VC{infiltration}: the syndicate to \V{broker} the new \V{hive} to. A tied \V{hivejack}: the \V{hive} that survives, or is \V{absorbed} next.}
\kr{Letters}{Buy \V{coin}: each letter one \V{coin} of that syndicate, up to \Rule{maxbuy} in any mix.}
\kr{\key{BACKSPACE}}{Take back the last \V{coin} in a buy.}
\kr{\key{RETURN}}{Confirm the buy. \key{RETURN} alone buys nothing.}
\kr{\key{Y}\ \ \key{N}}{When the game may end: \V{nullroute} the region now, or play on.}
\kr{The controller}{The stick moves the cursor: an \V{exploit} in your \V{kit} (a box on its \V{node}), an offered syndicate, or, while buying, a syndicate with a \V{hive}. The left fire button gives its answer, or YES; the right fire button takes back a \V{coin}, or NO; \key{\#} confirms.}
\end{keytable}

\subsection{When your \V{hive} is \V{absorbed}}
\begin{keytable}
\kr{Digits, \key{RETURN}}{First, how many \Vp{coin} to \V{cashout}; then how many to \V{swap}, an even number.
The rest you \V{hodl}. On the controller, the stick counts them (the \V{swap} two at a time, no higher than you
can) and the left fire button confirms.}
\end{keytable}

\subsection{Leaving}
\begin{keytable}
\kr{\key{ESCAPE}\ \ \key{Y}}{\VC{logoff} for good. A bot takes over your seat at once, and you're back on the title screen.
Any other key stays. On the controller, \key{*} with \key{\#}, then the left fire button; the right fire button stays.}
\end{keytable}

% ================================================================== 05
\renewcommand{\secstatus}{THE \MakeUppercase{\V{grid}} \quad // \quad THE SYNDICATES \quad // \quad SETUP}
\chapter{What's in the Game}
\begin{multicols}{2}
\rules{contents}
\end{multicols}

\subsection{The Seven Syndicates}
\syntable

\subsection{Setup}
\begin{multicols}{2}
\rules{setup}
\end{multicols}

% ================================================================== 06
\renewcommand{\secstatus}{\MakeUppercase{\V{spike}} \quad // \quad BUY \quad // \quad DRAW}
\chapter{Playing a Turn}
\begin{multicols}{2}
\rules{turn}
\end{multicols}

% ================================================================== 07
\renewcommand{\secstatus}{\MakeUppercase{\V{totalcompromise}} \quad // \quad \MakeUppercase{\V{lockdown}} \quad // \quad \MakeUppercase{\V{cashout}}}
\chapter{Ending the Game}
\begin{multicols}{2}
\rules{ending}
\end{multicols}

% ================================================================== 08
\renewcommand{\secstatus}{THE LOBBY \quad // \quad THE HOST \quad // \quad SEATS AND BOTS}
\chapter{Hosting, Joining and Leaving}
\begin{multicols}{2}
\rules{hosting}
\end{multicols}

% ================================================================== 09
\renewcommand{\secstatus}{\MakeUppercase{\V{finderscut}} \quad // \quad \MakeUppercase{\V{root}} AND \MakeUppercase{\V{sudo}} \quad // \quad CASH}
\chapter{Where the Money Comes From}
\begin{multicols}{2}
\rules{money}
\end{multicols}

% ================================================================== 10
\renewcommand{\secstatus}{EARLY \quad // \quad MIDDLE \quad // \quad LATE}
\vspace{24pt}\section{Tips}                  % follows on: the money chapter ends early on its page
\begin{multicols}{2}
\rules{tips}
\end{multicols}

% ================================================================== 11
\renewcommand{\secstatus}{THE RULES \quad // \quad THE ADAM}
\chapter{Frequently Asked Questions}
\begin{multicols}{2}
\rules{faq}

\subsection{On the Adam}
\begin{description}[leftmargin=0pt,labelsep=0.5em]
  \item[Which machines?] A Coleco Adam with a FujiNet for the Adam, and any TV. The shot clock counts the NTSC
    Adam's 60 frames a second; on a PAL Adam it counts a little slow, though the server keeps the real time.
  \item[Keyboard or controller?] Either, or both at once. The keyboard types the answers as the prompt line names
    them; the controller in port 1 moves a cursor over them.
  \item[Can I play alone?] Yes. Host a game and start it with \Rule{minplayers} seats: one bot plays you. Any seat
    nobody takes goes to a bot.
  \item[What are the digits on the \V{grid}?] Move hints: each marks a \V{node} that an \V{exploit} in your \V{kit}
    can run on, and the digit is its place in the \V{kit}. SmartKey \key{II} turns them on or off.
  \item[What is the yellow box on the \V{grid}?] The \V{node} the \V{exploit} under the cursor would \V{spike}.
  \item[Why are some numbers yellow?] Cash, and every other syndicate's \V{coin}, so two numbers side by side
    read apart: \code{3} next to \code{12} is not \code{312}.
  \item[What do \code{*} and \code{?} before a name mean?] \code{*} is a bot. \code{?} is a \V{player} who is away,
    or on autopilot; a bot plays for them meanwhile.
  \item[Can I turn the sound off?] Yes: SmartKey \key{III}, any time. A chord says it's back on.
  \item[Can I play against an Atari, an Apple II, a PC, a CoCo or an NES?] Yes. The servers are the same for every
    machine: all of them meet at the same tables.
  \item[Why was my handle refused?] Someone at the table has it, a bot included. If it's your own seat after a
    reset, the server may not have noticed the drop yet: try again in a minute.
  \item[What if everybody leaves?] The game waits for 10 minutes with nobody connected, then it is dropped.
\end{description}
\end{multicols}

% ================================================================== 12
\renewcommand{\secstatus}{EVERY TERM, FROM THE THEME}
\chapter{Glossary}
\rules{glossary}

% ================================================================== REF
\renewcommand{\secstatus}{FIND THE \MakeUppercase{\V{hive}}'S SIZE IN ITS TIER, THEN READ ACROSS}
\clearpage\renewcommand{\secnum}{REF}% after the page break: the running head uses it
\section{\V{coin} Prices and Payouts}
\vspace{6pt}
\begin{center}
\def\textcolor#1#2{\inv{#2}}% the tier's syndicate letters inverse, as on the screen
\newcommand{\hot}[1]{\leavevmode{\color{alarm}#1}}
\machine\fontsize{12}{14}\selectfont
\renewcommand{\arraystretch}{1.3}
\newcolumntype{C}{>{\centering\arraybackslash}p{0.85in}}
\newcolumntype{R}{>{\raggedleft\arraybackslash}p{0.95in}}
\begin{tabular}{CCC@{\hspace{14pt}}R@{\hspace{14pt}}RR}
LOW & MID & HIGH & \leavevmode\color{steel}PRICE & \leavevmode\color{signal}\MakeUppercase{\V{root}} & \leavevmode\color{signal}\MakeUppercase{\V{sudo}} \\
\TierCoins{low} & \TierCoins{mid} & \TierCoins{high} & & & \\[2pt]
\hline
2      & -      & -      & \Y{200}   & \Y{2,000}  & \Y{1,000} \\
3      & 2      & -      & \Y{300}   & \Y{3,000}  & \Y{1,500} \\
4      & 3      & 2      & \Y{400}   & \Y{4,000}  & \Y{2,000} \\
5      & 4      & 3      & \Y{500}   & \Y{5,000}  & \Y{2,500} \\
6-10   & 5      & 4      & \Y{600}   & \Y{6,000}  & \Y{3,000} \\
\hot{11-20}  & 6-10   & 5      & \Y{700}   & \Y{7,000}  & \Y{3,500} \\
21-30  & \hot{11-20}  & 6-10   & \Y{800}   & \Y{8,000}  & \Y{4,000} \\
31-40  & 21-30  & \hot{11-20}  & \Y{900}   & \Y{9,000}  & \Y{4,500} \\
41+    & 31-40  & 21-30  & \Y{1,000} & \Y{10,000} & \Y{5,000} \\
-      & 41+    & 31-40  & \Y{1,100} & \Y{11,000} & \Y{5,500} \\
-      & -      & 41+    & \Y{1,200} & \Y{12,000} & \Y{6,000} \\
\hline
\end{tabular}
\end{center}
\vspace{10pt}
\begin{center}\machine\fontsize{10}{16}\selectfont
\begin{tabular}{@{}ll@{}}
\textcolor{alarm}{RED} & FIRST \MakeUppercase{\V{safe}} SIZE: THE \MakeUppercase{\V{hive}} CAN'T BE \MakeUppercase{\V{absorbed}}\\
\textcolor{signal}{\MakeUppercase{\V{root}}} & 10 x \MakeUppercase{\V{coin}} PRICE\\
\textcolor{signal}{\MakeUppercase{\V{sudo}}} & 5 x \MakeUppercase{\V{coin}} PRICE\\
\Rule{endsize}+ & \MakeUppercase{\V{totalcompromise}}: THE \MakeUppercase{\V{nullroute}} MAY BE CALLED
\end{tabular}\end{center}

\vspace{6pt}
\begin{center}
{\machine\fontsize{9}{11}\selectfont\color{steel} THE LETTERS: EACH SYNDICATE ON THE \MakeUppercase{\V{grid}}, AND ITS TIER}\par\vspace{6pt}
\renewcommand{\arraystretch}{1.3}
\newcommand{\synrow}[1]{\inv{\Syn{#1}{letter}} & \Syn{#1}{name} & {\machine\color{steel}\MakeUppercase{\Syn{#1}{tier}}}}
\begin{tabular}{@{}cl@{\hspace{8pt}}l@{\hspace{36pt}}cl@{\hspace{8pt}}l@{}}
\synrow{0} & \synrow{4} \\
\synrow{1} & \synrow{5} \\
\synrow{2} & \synrow{6} \\
\synrow{3} & & & \\
\end{tabular}
\end{center}

\end{document}
